% =====================================================================
%  Приложение к ВКР: рекомендации по интеграции и практическому
%  применению разработанной методики в корпоративных сетях.
%
%  Сборка: xelatex -> biber -> xelatex -> xelatex
%  Преамбула общая со Stage1.
% =====================================================================
\documentclass[14pt, a4paper]{extarticle}

% =====================================================================
%  Подключаемые пакеты LaTeX
%  Сборка: XeLaTeX + Biber
%  Соответствие: ГОСТ 7.32-2017, ГОСТ Р 7.0.100-2018
% =====================================================================

% Кодировка и язык
\usepackage{fontspec}
\setmainfont{Times New Roman}
\setsansfont{Arial}
\setmonofont{Courier New}

\usepackage{polyglossia}
\setdefaultlanguage{russian}
\setotherlanguage{english}

% Геометрия страницы (ГОСТ 7.32-2017, п. 6.1):
% левое 30 мм, правое 15 мм, верхнее 20 мм, нижнее 20 мм
\usepackage[
    a4paper,
    left=30mm,
    right=15mm,
    top=20mm,
    bottom=20mm,
    footskip=10mm
]{geometry}

% Межстрочный интервал 1.5 (ГОСТ 7.32-2017, п. 6.3)
\usepackage{setspace}
\onehalfspacing

% Абзацный отступ 1.25 см
\usepackage{indentfirst}
\setlength{\parindent}{1.25cm}

% Математика
\usepackage{amsmath}
\usepackage{amssymb}
\usepackage{amsthm}
\usepackage{mathtools}

% Графика и таблицы
\usepackage{graphicx}
\usepackage{booktabs}
\usepackage{longtable}
\usepackage{tabularx}
\usepackage{array}
\usepackage{multirow}
\usepackage{makecell}
\usepackage{caption}
\usepackage{subcaption}
\usepackage{float}

% Цвет, ссылки, листинги
\usepackage{xcolor}
\usepackage[hidelinks]{hyperref}
\hypersetup{
    pdfencoding=auto,
    bookmarksnumbered=true,
    bookmarksopen=true,
    colorlinks=false
}
\usepackage{listings}
\lstset{
    basicstyle=\ttfamily\footnotesize,
    breaklines=true,
    showstringspaces=false,
    keywordstyle=\bfseries,
    columns=fullflexible,
    captionpos=b,
    frame=single,
    aboveskip=8pt,
    belowskip=8pt
}

% Перечисления и списки
\usepackage{enumitem}
\setlist{nosep, leftmargin=*}

% Заголовки и оглавление
\usepackage{titlesec}
\usepackage{tocloft}

% Библиография по ГОСТ Р 7.0.100-2018
\usepackage[
    backend=biber,
    style=gost-numeric,
    language=auto,
    sorting=none,
    maxbibnames=4,
    minbibnames=1
]{biblatex}

% Дополнительные служебные пакеты
\usepackage{csquotes}
\usepackage{etoolbox}
\usepackage{url}
\usepackage{microtype}
\usepackage{xurl}

% =====================================================================
%  Оформление по ГОСТ 7.32-2017
%  Структурные элементы: разделы, рисунки, таблицы, формулы, списки
% =====================================================================

% Размер основного шрифта 14 pt с межстрочным интервалом 1.5 устанавливается
% в classfile (extarticle/extreport, базовый размер 14pt). Здесь — стилевые
% настройки заголовков и нумерации.

% Заголовки по ГОСТ 7.32-2017 (п. 6.2):
%   - заголовки разделов и подразделов с абзацного отступа,
%     с прописной буквы, без точки в конце, не подчёркиваются;
%   - расстояние между заголовком и текстом — не менее одного интервала;
%   - заголовки разделов набираются жирным шрифтом, размер 14 pt.

\titleformat{\section}
    {\normalfont\fontsize{14}{17}\bfseries}
    {\thesection}{1em}{}
\titlespacing*{\section}{\parindent}{18pt}{12pt}

\titleformat{\subsection}
    {\normalfont\fontsize{14}{17}\bfseries}
    {\thesubsection}{1em}{}
\titlespacing*{\subsection}{\parindent}{14pt}{10pt}

\titleformat{\subsubsection}
    {\normalfont\fontsize{14}{17}\bfseries}
    {\thesubsubsection}{1em}{}
\titlespacing*{\subsubsection}{\parindent}{12pt}{8pt}

% Нумерация рисунков и таблиц — сквозная в пределах раздела
\renewcommand{\thefigure}{\arabic{section}.\arabic{figure}}
\renewcommand{\thetable}{\arabic{section}.\arabic{table}}
\renewcommand{\theequation}{\arabic{section}.\arabic{equation}}

% Подписи рисунков и таблиц по ГОСТ:
%   - «Рисунок N — Название» (тире), снизу, по центру;
%   - «Таблица N — Название», сверху, слева.
\captionsetup[figure]{
    labelsep=endash,
    name=Рисунок,
    justification=centering,
    singlelinecheck=true,
    font=normalsize
}
\captionsetup[table]{
    labelsep=endash,
    name=Таблица,
    justification=raggedright,
    singlelinecheck=false,
    font=normalsize
}

% Колонтитулы — простые, номер страницы по центру нижнего поля
\usepackage{fancyhdr}
\pagestyle{fancy}
\fancyhf{}
\fancyfoot[C]{\thepage}
\renewcommand{\headrulewidth}{0pt}
\renewcommand{\footrulewidth}{0pt}

% Оглавление: «СОДЕРЖАНИЕ» по центру, прописными
\renewcommand{\contentsname}{\hfill СОДЕРЖАНИЕ\hfill}

% Список литературы: «СПИСОК ИСПОЛЬЗОВАННЫХ ИСТОЧНИКОВ»
\defbibheading{bibliography}[Список использованных источников]{%
    \section*{#1}\addcontentsline{toc}{section}{#1}%
}

% =====================================================================
%  Пользовательские команды и сокращения
% =====================================================================

% Краткие ссылки
\newcommand{\figref}[1]{рисунок~\ref{#1}}
\newcommand{\Figref}[1]{Рисунок~\ref{#1}}
\newcommand{\tabref}[1]{таблица~\ref{#1}}
\newcommand{\Tabref}[1]{Таблица~\ref{#1}}
\newcommand{\eqnref}[1]{(\ref{#1})}
\newcommand{\secref}[1]{раздел~\ref{#1}}

% Английские сокращения как термины
\newcommand{\eng}[1]{\textit{#1}}

% Метки общих обозначений
\DeclareMathOperator*{\argmax}{arg\,max}
\DeclareMathOperator*{\argmin}{arg\,min}
\DeclareMathOperator{\softmax}{softmax}
\DeclareMathOperator{\sigm}{\sigma}
\DeclareMathOperator{\KL}{KL}
\DeclareMathOperator{\MMD}{MMD}
\DeclareMathOperator{\PSI}{PSI}

% Векторы и матрицы — обычным начертанием в соответствии с ГОСТ 7.32
% (без полужирного, как принято в большинстве отечественных НИР)

% Окружение для алгоритмов
\newtheoremstyle{algodef}{8pt}{8pt}{\normalfont}{}{\bfseries}{.}{0.5em}{}
\theoremstyle{algodef}
\newtheorem{algorithm}{Алгоритм}[section]
\newtheorem{definition}{Определение}[section]
\newtheorem{statement}{Утверждение}[section]


\addbibresource{../Stage1/bib/refs.bib}

\hypersetup{
    pdftitle={Рекомендации по интеграции и применению},
    pdfsubject={NIDS, SOC integration, NetFlow, IPFIX, Kafka, FastAPI},
}

\begin{document}

\begin{titlepage}
\centering
\vspace*{1cm}
{\large\bfseries Приложение к выпускной квалификационной работе}\\[0.5cm]
\rule{\linewidth}{0.4pt}\\[0.5cm]
{\Large\bfseries Рекомендации по интеграции и практическому применению
методики обнаружения несанкционированных действий в корпоративных
сетях}\\[0.5cm]
\rule{\linewidth}{0.4pt}\\[2cm]

\vfill
{\large \the\year}
\end{titlepage}

\tableofcontents
\newpage

\section*{ВВЕДЕНИЕ}
\addcontentsline{toc}{section}{ВВЕДЕНИЕ}

Документ содержит инженерные рекомендации по внедрению разработанной
методики обнаружения несанкционированных действий в продуктивную
инфраструктуру корпоративной сети. Рекомендации сформированы по
результатам экспериментов E1--E6 настоящей ВКР и согласуются с практикой
эксплуатации NTA-подсистем SOC.

% =====================================================================
\section{Архитектура развёртывания}

\subsection{Поточная схема}

В типовой корпоративной инфраструктуре прототип развёртывается между
сборщиками сетевой телеметрии и подсистемой алертирования SIEM/SOAR:

\begin{equation*}
\boxed{\text{Network sensors} \to \text{Collector} \to \text{Bus}
\to \text{Detector} \to \text{Alert sink}}
\end{equation*}

\begin{itemize}
    \item \textbf{Network sensors:} зеркальные SPAN-порты на ядре
    коммутации, сетевые TAP-устройства или host-based агенты типа
    \texttt{filebeat}/\texttt{nProbe}/\texttt{Zeek}.
    \item \textbf{Collector:} стандартные приёмники NetFlow/IPFIX
    (см. \cite{rfc3954, rfc7011}), напр. \texttt{nProbe Cento},
    \texttt{ntopng}, \texttt{Elastiflow}.
    \item \textbf{Bus:} потоковая очередь (\texttt{Apache Kafka},
    \texttt{NATS}, \texttt{Redis Streams}) для развязки сбора и анализа.
    \item \textbf{Detector:} FastAPI-сервис прототипа (\texttt{POST /score})
    в одном или нескольких репликах за \texttt{nginx}/HAProxy.
    \item \textbf{Alert sink:} SIEM (\texttt{Splunk, Elastic, QRadar})
    или SOAR-платформа, принимающая JSON-алерты по REST/Kafka.
\end{itemize}

\subsection{Минимальная конфигурация для пилота}

\begin{table}[H]
\centering
\caption{Минимальные требования к стенду пилотного развёртывания}
\renewcommand{\arraystretch}{1.2}
\begin{tabular}{ll}
\toprule
\textbf{Компонент} & \textbf{Минимум} \\
\midrule
CPU      & 4 ядра, AVX2 \\
RAM      & 8 GB \\
Диск     & 50 GB SSD (logs/checkpoints) \\
ОС       & Ubuntu 22.04 LTS / RHEL 9 / Windows 11 \\
Python   & 3.10--3.13 \\
GPU      & не обязательна; ускоряет обучение, не инференс \\
Сеть     & 1 Gbps (для приёма IPFIX) \\
\bottomrule
\end{tabular}
\end{table}

Эксперимент~E6 показал latency CNN-LSTM 0{,}96\,мс/окно на CPU при
batch\,$=$\,1; такая конфигурация поддерживает порядка тысяч
flow-в-секунду без батчирования и десятки тысяч с батчированием.

% =====================================================================
\section{Подключение источников трафика}

\subsection{NetFlow / IPFIX (рекомендуемый формат)}

Сборщик настраивается на экспорт всех 49 полей UNSW-NB15-схемы либо
её подмножества, общего с CICIDS2017 (12 признаков, см. Stage~2
§2.2). Минимальный набор полей, без которого инференс невозможен:

\begin{itemize}
    \item длительность потока (\emph{Flow Duration}),
    \item байты и пакеты в обе стороны
    (\emph{Total Length of Fwd/Bwd Packets}, \emph{Total Fwd/Bwd Packets}),
    \item межпакетные интервалы (\emph{Fwd/Bwd IAT Mean/Std}),
    \item TCP-флаги и состояние сессии (\emph{state}).
\end{itemize}

Маппинг IPFIX-полей в схему модели зафиксирован в
\verb|configs/data/unsw_nb15.yaml| и \verb|configs/data/cicids2017.yaml|.

\subsection{Сырые pcap-источники}

При наличии pcap-источников (например, по SPAN-портам) рекомендуется
извлекать flow-признаки утилитой \texttt{CICFlowMeter} либо
\texttt{Argus}: их выходные форматы напрямую совместимы с
\verb|configs/data/cicids2017.yaml|. Для UNSW-NB15-совместимого формата
доступен \texttt{nfdump}/\texttt{nfcapd} или скрипт извлечения,
поставляемый авторами датасета~\cite{moustafa2015unsw}.

\subsection{Endpoint-агенты}

При невозможности зеркалирования трафика на ядре можно использовать
host-based сборщики (\texttt{filebeat/IPFIX module},
\texttt{osquery + processes}, \texttt{Zeek conn.log}). Признаковый
вектор формируется на агенте или collector-стороне и отправляется в
Kafka-топик \texttt{nids.flows.v1}.

% =====================================================================
\section{Формирование и приоритизация алертов}

\subsection{Уровни severity}

Прототип реализует пятиступенчатую шкалу \emph{severity}, привязанную
к нормированному превышению скоринга над порогом:
\[
e = \frac{s(x) - \tau}{1 - \tau},
\quad
\text{severity} =
\begin{cases}
\text{info},     & s < \tau \\
\text{low},      & 0 \le e < 0{,}1 \\
\text{medium},   & 0{,}1 \le e < 0{,}3 \\
\text{high},     & 0{,}3 \le e < 0{,}6 \\
\text{critical}, & e \ge 0{,}6
\end{cases}
\]

Это позволяет аналитику SOC расставить приоритеты обработки без
дополнительных правил корреляции.

\subsection{Дедупликация и шумоподавление}

\texttt{AlertFormer} реализует временное окно дедупликации
(\(\Delta t = 5\)\,с по умолчанию): два алерта с одним и тем же
\emph{attack\_cat\_pred} и \emph{severity} в пределах окна
объединяются в один. На демо-прогоне (300 тиков, 8\,219 записей)
зафиксировано всего \textbf{30 алертов} --- эксплуатационно
приемлемая нагрузка.

\subsection{Структура алерта}

Формат JSON (одно событие = одна строка):

\begin{lstlisting}[basicstyle=\ttfamily\small]
{
    "alert_id": "alert-00000001",
    "ts": 1778468472.378,
    "severity": "medium",
    "score": 0.7102,
    "decision": 1,
    "fsm_state": "ATTACK_DDOS",
    "attack_cat_pred": "DoS",
    "profile": { "tick": 60, "...": "..." }
}
\end{lstlisting}

Поле \texttt{profile} переносится в SIEM как нативный объект; ключи
позволяют автоматически связать алерт с источником и тип-таксономией
SOC (например, MITRE ATT\&CK\cite{mitreattack}).

% =====================================================================
\section{Калибровка порогов и переобучение}

\subsection{Целевой режим эксплуатации}

В E3 порог \(\tau\) откалиброван на валидационной выборке под
целевой FPR\,$=$\,0{,}01. Для конкретной инфраструктуры рекомендуется:

\begin{enumerate}
    \item Зафиксировать допустимый объём FP-в-сутки на основе
    производительности команды SOC (типовое значение:
    \(\leq 50\) тяжёлых алертов на одного аналитика в смену).
    \item Перевести его в FP-per-time, разделив на средний поток
    flow-в-сутки коллектора.
    \item Использовать \texttt{find\_threshold\_for\_target\_fpr} на
    производственной валидационной выборке для подбора \(\tau\).
    \item Применять температурное масштабирование
    (\texttt{TemperatureScaler}) для согласования \(\hat{p}(y|x)\) с
    эмпирической вероятностью.
\end{enumerate}

\subsection{Drift-monitoring и автоматическое переобучение}

Блок \texttt{drift\_report} в эксплуатации формирует одну метрику
PSI/MMD на каждые \(N\)\,тыс. flow. Рекомендованные действия по уровням:

\begin{itemize}
    \item \(\text{PSI} \le 0{,}1\): нет необходимости в действиях.
    \item \(0{,}1 < \text{PSI} \le 0{,}25\): включается расширенное
    логирование, аналитик уведомляется.
    \item \(\text{PSI} > 0{,}25\): запускается процедура
    переобучения на свежей валидационной выборке; до переобучения
    severity повышается на одну ступень (компенсация снижения PR-AUC).
\end{itemize}

\subsection{Циклы переобучения}

Регулярное переобучение --- раз в квартал на накопленной
размеченной выборке; экстренное --- по факту срабатывания
drift-monitor. Размеры выборок --- не менее train-партиции UNSW-NB15
(\(\sim 80\,000\) flows) для устойчивого качества.

% =====================================================================
\section{Безопасность и приватность}

\subsection{Принцип ``данные не покидают периметр''}

Прототип не передаёт сырые признаки во внешние сервисы. Все артефакты
(\texttt{preprocessor.json}, чекпойнты \texttt{models/*.pt}) хранятся
локально и шифруются на уровне диска (LUKS / BitLocker).

\subsection{Аутентификация REST-API}

В пилотном развёртывании рекомендуется поставить FastAPI за
reverse-proxy с mTLS либо использовать встроенную в FastAPI
авторизацию по API-token. Без аутентификации сервис не должен быть
доступен из внешних сегментов.

\subsection{Аудит}

Все алерты дублируются в \texttt{experiments/runs/realtime\_alerts.jsonl}
для офлайн-аудита. Drift-readings --- в
\texttt{realtime\_drift.jsonl}. Журналы выбрасываются в SIEM
без изменения.

% =====================================================================
\section{Развёртывание в Docker / Kubernetes}

\subsection{Docker Compose (рекомендуется для пилота)}

\begin{lstlisting}[basicstyle=\ttfamily\small]
version: "3.9"
services:
  detector:
    image: diploma-nids:0.1.0
    command: uvicorn diploma_nids.inference.service:app --host 0.0.0.0 --port 8000
    environment:
      DIPLOMA_MODEL_YAML: configs/models/cnn_lstm.yaml
      DIPLOMA_CHECKPOINT: models/cnn_lstm_seed42.pt
      DIPLOMA_PREPROCESSOR: data/processed/unsw/preprocessor.json
      DIPLOMA_THRESHOLD: "0.675"
    volumes:
      - ./models:/app/models:ro
      - ./data/processed:/app/data/processed:ro
      - ./experiments/runs:/app/experiments/runs
    ports: [ "8000:8000" ]

  ui:
    image: diploma-nids:0.1.0
    command: streamlit run src/diploma_nids/ui/streamlit_app.py --server.port 8501 --server.address 0.0.0.0
    environment:
      DIPLOMA_API_URL: http://detector:8000
    ports: [ "8501:8501" ]
    depends_on: [ detector ]
\end{lstlisting}

\subsection{Масштабирование}

Поскольку инференс stateless (за исключением \texttt{WindowScorer},
который держит rolling buffer per-client), сервис горизонтально
масштабируется. Sharding по \texttt{src\_ip} либо по \texttt{client\_id}
обеспечивает корректность буферов.

% =====================================================================
\section{Воспроизводимость экспериментов}

Полный воспроизводимый цикл одной командой:

\begin{lstlisting}[basicstyle=\ttfamily\small, language=bash]
make install-dev
python scripts/02_build_dataset.py \
    --config configs/data/unsw_nb15.yaml \
    --preprocess configs/preprocess/default.yaml \
    --strategy official
python scripts/10_run_experiment_E1.py \
    --train configs/train/full.yaml --seeds 42 123 2024
python scripts/11_run_experiment_E2_error_analysis.py
python scripts/12_run_experiment_E4_attacker.py
python scripts/13_run_experiment_E5_drift.py
python scripts/20_build_thesis_tables.py
python scripts/22_parse_training_logs.py
python scripts/21_plot_training_history.py
python scripts/30_demo_run.py
\end{lstlisting}

Все артефакты сохраняются в \texttt{results/} и \texttt{experiments/};
seeds фиксированы и согласованы между скриптами.

% =====================================================================
\section{Ограничения и направления развития}

\subsection{Ограничения текущего прототипа}

\begin{itemize}
    \item Модель обучена на UNSW-NB15 (\(\sim 70\,\text{тыс.}\) flow);
    для продуктивной среды требуется fine-tuning на свежем трафике.
    \item Шаблоны ИИ-злоумышленника находятся в домене,
    отличающемся от UNSW по маргинальным распределениям (E5 показал
    PSI > 0{,}25 во всех точках). Необходима калибровка шаблонов
    под целевую инфраструктуру.
    \item TLS-инспекция и работа на пакетном уровне не реализованы;
    методика ориентирована на flow-уровень.
\end{itemize}

\subsection{Направления развития}

\begin{itemize}
    \item \textbf{cGAN-составляющая агента}: замена части шаблонов
    обученной cGAN-моделью для более реалистичного синтетического
    трафика;
    \item \textbf{Cross-evaluation на CICIDS2017} в полном объёме с
    дообучением на 10\% выборки CICIDS;
    \item \textbf{Расширение на pcap-уровень} с использованием полного
    UNSW-NB15\_1..4.csv (\(\sim 2{,}5\,\text{млн}\) записей);
    \item \textbf{Поддержка дополнительных датасетов} (TON\_IoT,
    Bot-IoT, LITNET-2020) для проверки переносимости;
    \item \textbf{Online-обучение} с использованием
    EWC/replay-buffer для адаптации к concept drift без полного
    переобучения.
\end{itemize}

% =====================================================================
\printbibliography[title={СПИСОК ИСПОЛЬЗОВАННЫХ ИСТОЧНИКОВ}]

\end{document}
